\subsection{Tuning the Connected-Component Recognition Approach}
\label{sec:cca-search}

The final recognition experiment investigated whether spatial measurements
could improve the tuned colour classifier. The preceding search selected
configuration 2153, which correctly classified 81/86 matched ball types and
74/86 numbers on Table1--Table6. Rather than replacing this classifier with
an entirely different method, CCA was evaluated as an optional type-correction
stage. Its purpose was to distinguish a substantial white stripe from smaller
white markings and highlights using connectivity and boundary contact.

The earlier default CCA method reduced accuracy, but that result only tested
one set of thresholds. This experiment separately investigated feature
extraction, decision thresholds and which correction rules should be enabled.
All changes remained in the independent experiment executable. The production
CCA implementation and main application were not modified.

\subsubsection{Fixed recognition base and control configurations}

The primary search froze configuration 2153: K=6; light-cluster thresholds
$V\geq170$, $S\leq125$; an 85\% cue-ball light threshold; an 80\%
solid-ball non-light threshold; and a blue/purple hue boundary of 125.
The other colour thresholds retained their original values. Extraction
used the original Hough radius with zero padding and no cloth filtering.
Neither the colour clusters nor body-cluster selection nor number thresholds
were tuned again. CCA could change the type passed to number recognition,
but could not change the cluster colours themselves.

Four controls were measured: the original K=2 classifier without CCA; that
classifier with production CCA; the tuned K=6 classifier without CCA; and
the tuned K=6 classifier with the original CCA thresholds. For the last
control, the definition of a light cluster was aligned with the K=6 base
($S\leq125$). It was not left at the original K=2 saturation threshold
of 165. This alignment avoids constructing CCA features from a different
light/non-light partition than the classifier whose decisions they modify.

Detection and one-to-one annotation matching were performed once per image
and reused for all CCA settings. The run again contained 86 matched
annotations out of 90 on Table1--Table6, plus 16 matched annotations on
AllBalls. Unmatched detections were displayed but excluded from classification
accuracy. Table1--Table6 provided tuning scores; AllBalls remained a separate
reference check and did not influence ranking. The same limitations of
correlated scenes and reuse of evaluation labels described in
Section~\ref{sec:recognition-experiments} apply here.

\subsubsection{Constructing and measuring connected regions}

K-means clusters contain pixels with similar BGR colours, regardless of their
positions. CCA introduces spatial information after clustering. All cluster
masks classified as light were combined with a binary OR operation, then
restricted to the ball mask. The non-light mask was the complement of the
light mask within the ball. Adjacent light pixels from different K-means
clusters could therefore form one connected light region. Conversely,
disconnected patches from one colour cluster could form several regions.

Eight-connected component labelling was applied separately to the light and
non-light masks. A component's area was divided by the retained ball-mask
area, rather than by the area of its own colour category. The features were
the number of eligible light components, the largest eligible light area,
the number of eligible non-light components and the largest eligible
non-light area. Eligibility required a minimum component area of
$\{0.5,1,2,3\}\%$ of the ball.

The circular mask was eroded with an elliptical kernel, and the eroded mask
was subtracted from the original mask to define an inner boundary band.
Kernel sizes $3\times3$, $5\times5$ and $7\times7$ compared successively
wider bands. The original erosion-border semantics were retained. Light
boundary coverage was the percentage of band pixels belonging to the raw
light mask:
\[
 C_{\partial}=100\frac{|M_{\mathrm{light}}\cap B|}{|B|},\qquad
 B=M_{\mathrm{ball}}\setminus\operatorname{erode}(M_{\mathrm{ball}}).
\]
An empty band returned zero coverage.

Small-component filtering affected component counts and largest-component
statistics only. It did not erase pixels from the light/non-light masks,
change their raw percentages, or change boundary coverage. This distinction
preserved the original feature definitions and prevented an area-cutoff
change from silently altering several measurements at once. No morphological
opening or closing was added to the colour masks.

The four area cutoffs and three kernels gave 12 feature configurations.
These were cached for each candidate after K-means, so rule-threshold trials
reused the same spatial measurements. For the original K=2 control, the
experimental defaults were compared with production CCA on every candidate:
light and non-light masks, boundary mask, component counts, largest areas,
boundary coverage and final hybrid type all had to agree. This guarded
against improvements caused by accidentally changing the baseline algorithm.

\subsubsection{Decision rules and search space}

The hybrid retained the original decision order: test a cue-ball correction,
then a striped-ball correction, then a solid-ball correction, and otherwise
retain the colour-only type. Each rule could be enabled independently,
giving all eight subsets of the three rules, including CCA disabled.
This comparison distinguished improvements from a useful subset of rules
from improvements requiring the full hybrid.

The following threshold groups were searched:

\begin{enumerate}
\item \textbf{Cue rule.} Minimum raw light percentage
$\{80,85,90,95\}$; minimum largest light component
$\{70,80,90\}\%$ of the ball; and minimum light-boundary coverage
$\{60,75,90\}\%$. The rule also required zero non-light components
above the selected area cutoff. These conditions gave 36 combinations.
The rule is an override when its conditions hold, not a veto on every
other prediction of white: if no override fires, the base type remains.
\item \textbf{Stripe rule.} Minimum largest light component
$\{5,10,15,20,25,30\}\%$; minimum light-boundary coverage
$\{5,10,20,30,40,50\}\%$; and minimum total non-light proportion
$\{2,5,10,15,20\}\%$. These gave 180 combinations. The non-light
requirement prevents an almost entirely light ball from being automatically
assigned a stripe merely because its light region is connected to the edge.
\item \textbf{Solid rule.} Minimum total non-light proportion
$\{60,70,80,90\}\%$; maximum light-boundary coverage
$\{5,12,20,30\}\%$; maximum eligible non-light component count
$\{1,2\}$, with at least one required; and minimum largest non-light
component $\{0,50,70,85\}\%$. These gave 128 combinations.
A zero largest-component threshold disabled this additional dominance
constraint. Allowing two components and adding a dominance threshold were
experimental extensions; the original rule required exactly one component
and had no explicit largest-area condition.
\end{enumerate}

For every feature setting and enabled-rule subset, the original rule
thresholds supplied the starting point. Two passes searched the cue, stripe
and solid groups in that order, skipping disabled groups. Each group used
its full Cartesian product while holding the other groups fixed, and the
current configuration remained a candidate. This was a grouped search,
not a full joint Cartesian product over all rule thresholds. It may miss
interactions requiring simultaneous changes in several groups, and two
passes do not guarantee convergence to a global optimum.

For each rule, four of the eight subsets enable it. Consequently the run
scored
\[
 4 + 12\left[8+2\times4(36+180+128)\right]=33{,}124
\]
trials: four controls, 96 initial seeds, and 33,024 grouped threshold trials.
Repeated parameter settings remained in the audit trail. This total is
therefore not a count of distinct models. No thresholds were selected
separately for individual images or ball numbers.

Ranking first maximised correct full-number predictions, then correct type
predictions. Remaining ties preferred fewer number regressions, fewer type
regressions, fewer type changes, fewer enabled rules, less departure from
the original CCA parameters, and finally the earlier trial ID. All
regressions and changes were measured against configuration 2153. A
type-first ranking reversed the first two criteria. The departure metric
was only a deterministic tie-break: it summed absolute feature and threshold
changes using the fixed normalisations documented in
\code{experiments/ccaGridSearch.cpp}. It was not an additional measured
recognition score or a statistical regularisation estimate.

\subsubsection{Results and rule ablations}

% Generated by experiments/writeRecognitionReport.py from saved CSVs.
\begin{table}[htbp]
\centering
\small
\caption{CCA controls and selected configuration. Tuning scores use the 86 matched balls in Table1--6; the reference column reports numbers on AllBalls. The K=6 default-CCA control retains original CCA rule thresholds but aligns its light definition with configuration 2153.}
\label{tab:cca-search-summary}
\begin{tabular}{lrrrrr}
\toprule
Configuration & Type & Type acc. & Number & Number acc. & Ref. no. \\
\midrule
Original K=2 & 76/86 & 88.37\% & 67/86 & 77.91\% & 16/16 \\
Original K=2 + CCA & 73/86 & 84.88\% & 64/86 & 74.42\% & 16/16 \\
Tuned K=6 & 81/86 & 94.19\% & 74/86 & 86.05\% & 16/16 \\
Tuned K=6 + default CCA & 79/86 & 91.86\% & 72/86 & 83.72\% & 16/16 \\
Tuned K=6 + selected CCA & 82/86 & 95.35\% & 75/86 & 87.21\% & 16/16 \\
\bottomrule
\end{tabular}
\end{table}
\FloatBarrier

% Generated by experiments/writeRecognitionReport.py from saved CSVs.
\begin{table}[htbp]
\centering
\small
\caption{Best observed configuration for each enabled-rule subset on the K=6 base. Fixes and regressions compare with configuration 2153. Each subset is tuned separately.}
\label{tab:cca-search-ablation}
\begin{tabular}{lrrrrr}
\toprule
Rules enabled & Type & Number & Type fixes/losses & No. fixes/losses & ID \\
\midrule
None & 81 & 74 & 0/0 & 0/0 & 3 \\
Cue & 81 & 74 & 0/0 & 0/0 & 19326 \\
Stripe & 81 & 74 & 0/0 & 0/0 & 19503 \\
Cue + stripe & 81 & 74 & 0/0 & 0/0 & 19900 \\
Solid & 82 & 75 & 1/0 & 1/0 & 20242 \\
Cue + solid & 82 & 75 & 1/0 & 1/0 & 20535 \\
Stripe + solid & 82 & 75 & 1/0 & 1/0 & 21008 \\
Cue + stripe + solid & 82 & 75 & 1/0 & 1/0 & 21661 \\
\bottomrule
\end{tabular}
\end{table}
\FloatBarrier


With original K=2 recognition, default CCA reduced correct types from 76 to
73 and correct numbers from 67 to 64. With the K=6 light definition aligned
to configuration 2153, the default rules still reduced type correctness
from 81 to 79 and number correctness from 74 to 72. Thus simply attaching
the existing hybrid to the improved colour classifier was not beneficial.

After tuning, the best observed configuration correctly classified 82/86
types ($95.35\%$) and 75/86 numbers ($87.21\%$). Both number-first and
type-first selection chose configuration 20242. Relative to the tuned K=6
base, it corrected one type and its dependent number, with no regressions
on matched Table1--Table6 candidates. Each accuracy increased by 1.16
percentage points. End-to-end number recall increased from 74/90
($82.22\%$) to 75/90 ($83.33\%$), with detection unchanged. AllBalls
remained at 16/16 correct types and numbers.

The cue-only and stripe-only ablations did not improve on the K=6 base.
The solid-only rule achieved the best scores. Some larger tuned rule
subsets achieved the same scores, but did not provide additional correct
predictions. The tie-break favouring fewer enabled rules selected the
solid-only model. This is evidence for a limited corrective use of the
spatial feature, not for enabling every CCA override.

% Generated by experiments/writeRecognitionReport.py from saved CSVs.
\begin{table}[htbp]
\centering
\small
\caption{Original CCA settings and selected CCA settings on the fixed K=6 base. For enable flags, 1 means enabled and 0 means disabled.}
\label{tab:cca-search-parameters}
\begin{tabular}{lrr}
\toprule
Parameter & Default CCA & Selected CCA \\
\midrule
Minimum component area (\%) & 2.000000 & 2.000000 \\
Erosion kernel size & 5 & 5 \\
Cue rule enabled & 1 & 0 \\
Stripe rule enabled & 1 & 0 \\
Solid rule enabled & 1 & 1 \\
Solid minimum non-light (\%) & 70 & 70 \\
Solid maximum light boundary (\%) & 12 & 20 \\
Maximum non-light components & 1 & 1 \\
Minimum largest non-light region (\%) & 0 & 0 \\
\bottomrule
\end{tabular}
\end{table}
\FloatBarrier


The winning feature settings retained the original 2\% component cutoff
and $5\times5$ erosion kernel. The cue and stripe overrides were disabled.
The solid correction required at least 70\% non-light pixels, at most
20\% light-boundary coverage, and exactly one non-light component above
the area cutoff. No additional largest-non-light dominance threshold was
needed. The decisive threshold change relative to the original solid
rule was raising its allowable light-boundary coverage from 12\% to 20\%.
These correction thresholds are separate from the base classifier's 80\%
non-light threshold.

\subsubsection{Interpreting the corrected ball}

% Generated by experiments/writeRecognitionReport.py from saved CSVs.
\begin{table}[htbp]
\centering
\small
\caption{Matched candidates changed by the selected CCA configuration. Candidate IDs refer to detections rather than annotated ball numbers.}
\label{tab:cca-search-changes}
\begin{tabular}{lrrrrll}
\toprule
Image & Candidate & True no. & Base no. & CCA no. & Base type & CCA type \\
\midrule
Table1.jpeg & 10 & 5 & 13 & 5 & Striped & Solid \\
\bottomrule
\end{tabular}
\end{table}
\FloatBarrier


Only candidate 10 in Table1 changed under the selected correction. Its
annotated number was 5, a solid orange ball. The K=6 colour classifier
measured 21.83\% light and 78.17\% non-light pixels. Since the non-light
proportion fell below the base solid threshold of 80\%, it predicted
striped, causing the orange body colour to be assigned number 13.

CCA found three light components above the 2\% area cutoff, with the
largest occupying 12.74\% of the ball. It found one eligible non-light
component, occupying 76.11\% of the ball. Light pixels covered 17.68\%
of the boundary band. This failed the old solid correction's maximum
boundary coverage of 12\%, but satisfied the tuned maximum of 20\%.
Together with the 70\% non-light requirement and single-component
condition, this changed the type to solid and the dependent number to 5.

The interpretation is that a substantial connected body region and
relatively little light at the boundary provided useful evidence for a
solid ball, even though its total light proportion was too high for the
colour-only rule. This is a feature-based explanation of this candidate,
not proof that every similar mask corresponds to a solid ball.

The diagnostic crop also contains substantial parts of neighbouring balls
in the tightly arranged Table1 scene. Its connected non-light region cannot
therefore be assumed to represent only the target ball's body. A correct
classification does not establish an accurate segmentation. In this case,
the tuned rule may partly compensate for contamination from imperfect
localisation rather than isolate a property of an unobscured solid ball.

\begin{figure}[htbp]
\centering
\includegraphics[width=\linewidth]{cca_search_example.png}
\caption{The corrected Table1 candidate under the original K=2 control,
original CCA, tuned K=6, default CCA on K=6, and the selected solid-only
CCA correction. Each panel shows the masked ball and its light/non-light
partition. White and orange represent interior light and non-light pixels;
green and red in the mask represent light and non-light boundary pixels.
Prediction text uses green for correct and red for incorrect. The mask
colours are feature visualisations, not ground-truth segmentation labels.}
\label{fig:cca-search-example}
\end{figure}
\FloatBarrier

\subsubsection{Ceiling for type-only corrections and limitations}

The preceding colour experiment achieved 75/86 oracle-number predictions
when supplied with true types. The CCA experiment independently recomputed
the same diagnostic with its fixed K=6 features and number rules. For
these rules, a correct predicted number implies the correct type: 0
corresponds to white, 1--8 to solid, and 9--15 to striped. Consequently,
75/86 is an upper bound on full-number correctness achievable by changing
only the type of these cached candidates. The selected CCA correction
reached that measured ceiling. Further type-only changes cannot improve
full-number correctness on this fixed feature set; the remaining eleven
number errors require changes to colour interpretation, body selection,
segmentation or extraction. This is not an upper bound on a different
recognition pipeline, and it does not imply that all type errors have
been corrected: four type errors remain.

The observed gain is small and entirely due to one candidate in one
image. Searching over 33,124 trials on the same small tuning set creates
a substantial opportunity to select a threshold that fits this example.
The absence of regressions on these images is useful for comparison, but
does not establish that the rule will be harmless on new ball orientations
or lighting. An independent image set would be needed before treating
the one-ball gain as a general improvement. The report therefore records
the result as a promising targeted correction rather than a definitive
replacement for the colour-only classifier.

\subsubsection{Reproducing and inspecting the CCA experiment}

After building the independent experiment target, run from the repository root:
\begin{quote}
\code{./build-local/BallExperiments --cca-grid}
\end{quote}
An optional directory argument preserves separate runs. Feature ranges, rule thresholds and the
frozen recognition base are defined in \code{experiments/ccaSearchConfig.h}.
Outputs are written to \code{output/experiments/cca_grid_search/}.

The leaderboard records every trial's full CCA configuration, stage, parent,
scores and regression counts. The ablation CSV selects the best observed
model for each enabled-rule subset, and the seed CSV records each grouped
search's starting and final IDs. Finalist configurations, per-image metrics,
candidate predictions, confusion matrices, detection counts and cached
features support inspection of the results. The number-ceiling CSV records
the true-type diagnostic per image and per split. One diagnostic image per
candidate compares distinct control/finalist configurations, with no
individual ball subdirectories. Unmatched candidates have blank correctness
fields and unavailable correctness indicators.

The report tables and figure are refreshed with
\code{python3 experiments/writeCCAReport.py}. They are snapshots rather
than live links to the output directory. The surrounding prose contains
run-specific counts and should be reviewed after changing the grid. No
search result was automatically applied to the main application.
